\section{Bipolar Junction Transistor (BJT)}

A BJT is three-port, active, non-linear and time-variant device.
It is physically built from two PN junctions that share a single
region, with the two junctions oriented oppositely. This shows
up in its governing model. It is called a \emph{bipolar} device
because both electrons and holes are involved in the conduction
process.

The three terminals are the base (B), the collector (C) and the
emitter (E). It is often called a \emph{current controlled device}
because of its base-collector current relationship. In reality, the
base-emitter voltage controls the collector current.

BJT family has two types of transistors: NPN and PNP.\ They are
(in a sense) mirror images of each other.

\begin{circuitfig}
	% Paths, nodes and wires:
	\node[npn] at (5.5, 7.02){};
	\node[pnp, yscale=-1] at (9, 6.98){};
	\node[shape=rectangle, minimum width=1.465cm, minimum height=0.465cm] at (5.5, 6){} node[anchor=north, align=center, text width=1.077cm, inner sep=6pt] at (5.5, 6.25){Emitter};
	\node[shape=rectangle, minimum width=0.97cm, minimum height=0.465cm] at (4.157, 7.02){} node[anchor=north, align=center, text width=0.582cm, inner sep=6pt] at (4.157, 7.27){Base};
	\node[shape=rectangle, minimum width=0.97cm, minimum height=0.465cm] at (7.657, 6.98){} node[anchor=north, align=center, text width=0.582cm, inner sep=6pt] at (7.657, 7.23){Base};
	\node[shape=rectangle, minimum width=1.465cm, minimum height=0.465cm] at (5.5, 8.04){} node[anchor=north, align=center, text width=1.077cm, inner sep=6pt] at (5.5, 8.29){Collector};
	\node[shape=rectangle, minimum width=1.465cm, minimum height=0.465cm] at (9, 8){} node[anchor=north, align=center, text width=1.077cm, inner sep=6pt] at (9, 8.25){Collector};
	\node[shape=rectangle, minimum width=1.465cm, minimum height=0.465cm] at (9, 5.96){} node[anchor=north, align=center, text width=1.077cm, inner sep=6pt] at (9, 6.21){Emitter};
	\node[shape=rectangle, minimum width=1.965cm, minimum height=0.755cm] at (5.5, 4.855){} node[anchor=north, align=center, text width=1.577cm, inner sep=6pt] at (5.5, 5.25){\Large NPN};
	\node[shape=rectangle, minimum width=1.965cm, minimum height=0.755cm] at (9, 4.855){} node[anchor=north, align=center, text width=1.577cm, inner sep=6pt] at (9, 5.25){\Large PNP};
\end{circuitfig}

The behaviour of a BJT can be approximated well by the below equations.
\begin{itemize}
  \item \textbf{Mathematical Model}:
    There are several mathematical models that describe the behaviour of the
    BJT. Some of the most common ones are:
    \begin{itemize}
      \item \textbf{Ideal Model}: The ideal model assumes that the BJT has infinite
        current gain and zero base current.
        \[
          I_C = \begin{cases}
            0 & \text{if } V_{BE} < V_{th} \\
            \infty & \text{if } V_{BE} > V_{th}
          \end{cases}
        \]
      \item \textbf{Non-linear Model}: The Ebers-Moll model combined with Base-Width Modulation
        Non-Linearity is a more accurate model that describes the behaviour of BJT. Ebers-Moll
        is based on from two coupled Shockley's diode equations.
        \[
          I_C = \left[\underbrace{\alpha_F I_{FS} \left( e^{\frac{V_{BE}}{V_T}} - 1 \right)}_{\text{Forward Contribution}} - \underbrace{\alpha_R I_{RS} \left( e^{\frac{V_{BC}}{V_T}} - 1 \right)}_{\text{Reverse Contribution}}\right] \left( 1 + \frac{V_{CE}}{V_A} \right)
        \]
        \begin{itemize}
          \item \(I_C\) is the collector current
          \item \(I_{FS}\) is the forward saturation current
          \item \(\alpha_F\) is the forward common base current gain
          \item \(V_{BE}\) is the base-emitter voltage
          \item \(V_T\) is the thermal voltage
          \item \(I_{RS}\) is the reverse saturation current
          \item \(\alpha_R\) is the reverse common base current gain
          \item \(V_{BC}\) is the base-collector voltage
          \item \(V_{CE}\) is the collector-emitter voltage
          \item \(V_A\) is the Early voltage
        \end{itemize}
        We \textbf{will make two simplifications to our Non-Linear model} to make
        it easier to work with,
        \begin{itemize}
          \item Ignore the dependence of collector current on collector-emitter
            voltage (\(V_{CE}\)) by assuming an infinite early voltage \(V_A \to \infty\).
            This simplifies the Ebers-Moll model to:
            \[
              I_C = \alpha_F I_{FS} \left( e^{\frac{V_{BE}}{V_T}} - 1 \right) - \alpha_R I_{RS} \left( e^{\frac{V_{BC}}{V_T}} - 1 \right)
            \]
          \item In forward active mode, the base-collector voltage is reverse biased
            (\(V_{BC} \ll 0\)) and thus \( e^{\frac{V_{BC}}{V_T}} \approx 0 \) and the
            Ebers-Moll model simplifies to:
            \begin{align*}
              I_C &= \alpha_F I_{FS} \left( e^{\frac{V_{BE}}{V_T}} - 1 \right) - \alpha_R I_{RS} \left( e^{\frac{V_{BC}}{V_T}} - 1 \right) \\
              I_C &= \alpha_F I_{FS} \left( e^{\frac{V_{BE}}{V_T}} - 1 \right) - \alpha_R I_{RS}(-1) \\
              I_C &= \alpha_F I_{FS} \left( e^{\frac{V_{BE}}{V_T}} - 1 \right) + \alpha_R I_{RS} \\
              I_C &\approx \alpha_F I_{FS} \left( e^{\frac{V_{BE}}{V_T}} - 1 \right)
            \end{align*}
          \item Forward and Reverse Saturation Currents \( I_{FS} \) and \(I_{RS}\)
            are related to each other by the saturation current \(I_{S}\):
            \[
              \alpha_F I_{FS} = \alpha_R I_{RS} = I_S
            \]
            Thus, the collector current model becomes:
            \[
              I_C \approx I_S \left( e^{\frac{V_{BE}}{V_T}} - 1 \right)
            \]
        \end{itemize}
        We \textbf{can also derive the base and emitter currents} from the collector
        current model.
        \begin{itemize}
          \item Base Current (\(I_B\)):
            The base current is given by:
            \[
              I_B = \frac{I_C}{\beta}
            \]
            where \(\beta\) is the current gain of the BJT.
          \item Emitter Current (\(I_E\)):
            The emitter current is given by:
            \begin{align*}
              I_E &= I_C + I_B \\
              I_E &= I_C + \frac{I_C}{\beta} \\
              I_E &= I_C \left( 1 + \frac{1}{\beta} \right)
              I_E \approx I_C \hspace{1cm} \text{[for large \(\beta\)]}
            \end{align*}
        \end{itemize}
      \item \textbf{Linear Model}: The Hybrid-Pi model models the BJT as a
        dependent current source when linearized around a bias point to set
        its transconductance \(g_m\) property.
        \begin{align*}
          I_C = g_m V_{BE} + g_o V_{CE} \\
          I_B = g_\pi V_{BE}
        \end{align*}
        Here, \(g_m\) is the transconductance, \(g_o\) is the
        output conductance and \(g_\pi\) is the input conductance. Below
        schematic shows the simple idea of the hybrid-pi model for BJT.
        \begin{circuitfig}
        	\draw (9.25, 6.75) to[american current source, l_={$g_m V_{BE}$}] (9.25, 5);
        	\draw (6.5, 6.75) to[european resistor, l={$g_\pi$}] (6.5, 5);
        	\draw (10.75, 6.75) to[european resistor, l={$g_o$}] (10.75, 5);
        	\draw (6.5, 5) -- (9.25, 5);
        	\draw (9.25, 5) -- (10.75, 5);
        	\draw (10.75, 6.75) -- (9.25, 6.75);
        	\draw (10.75, 6.75) -- (11.75, 6.75);
        	\draw (6.5, 6.75) -- (5.75, 6.75);
        	\node[shape=rectangle, minimum width=0.965cm, minimum height=0.465cm](N1) at (5.25, 6.75){} node[anchor=center] at (N1.text){$V_B$};
        	\node[shape=rectangle, minimum width=0.965cm, minimum height=0.465cm](N2) at (8.5, 4.5){} node[anchor=center] at (N2.text){$V_E$};
        	\draw (8.5, 5) -| (8.5, 4.75);
        	\node[shape=rectangle, minimum width=0.965cm, minimum height=0.465cm](N3) at (12.25, 6.75){} node[anchor=center] at (N3.text){$V_C$};
        \end{circuitfig}
        \(g_m\) can be calculated as follows:
        \begin{align*}
          g_m &= \frac{dI_C}{dV_{BE}} |_{V_{DS(Q)}} \\
          g_m &= \frac{d}{dV_{BE}} \left[ I_S \left( e^{\frac{V_{BE}}{V_T}} - 1 \right) \right] |_{V_{DS(Q)}} \\
          g_m &= \frac{d}{dV_{BE}} \left[ I_S e^{\frac{V_{BE}}{V_T}} - I_S \right] |_{V_{DS(Q)}} \\
          g_m &= \left[ I_S e^{\frac{V_{BE}}{V_T}} \times \frac{1}{V_T} \right] \\
          g_m &= \frac{I_C}{V_T} \hspace{1cm}
        \end{align*}
        \(g_o\) can be calculated as follows:
        \begin{align*}
          g_o &= \frac{dI_C}{dV_{CE}} |_{V_{BE(Q)}} \\
          g_o &= \frac{d}{dV_{CE}} \left[ I_S \left( e^{\frac{V_{BE}}{V_T}} - 1 \right) \left( 1 + \frac{V_{CE}}{V_A} \right) \right] |_{V_{BE(Q)}} \\
          g_o &= \left[ I_S \left(e^{\frac{V_{BE}}{V_T}} - 1 \right) \times \frac{1}{V_A} \right] \\
          g_o &= \frac{I_C}{V_A} \hspace{1cm}
        \end{align*}
    \end{itemize}
  \item \textbf{Voltages \& Relations}:
    There are several important voltages in a BJT:\
    \begin{itemize}
      \item \textbf{Thermal/Boltzmann Voltage (\(V_T\))}:
        The thermal voltage is given by:
        \[
          V_T = \frac{kT}{q}
        \]
        where \(k\) is the Boltzmann constant, \(T\) is the temperature in
        Kelvin and \(q\) is the charge of an electron. At room temperature,
        \(V_T \approx 26\,\mathrm{mV}\):
      \item \textbf{Threshold Voltage (\(V_{th}\))}:
        The threshold voltage is the minimum voltage required to bias a junction
        such that it conducts. For a silicon junction (most common), the
        threshold voltage is \(V_{th} \approx 0.6\text{--}0.7\,\mathrm{V}\).
        Since we have two junctions in a BJT, we have two threshold voltages.
        We define \(V_{Fth}\) for the forward junction threshold voltage and
        \(V_{Rth}\) for the reverse junction threshold voltage.
    \end{itemize}
    They are also related to each other:
    \begin{itemize}
      \item \textbf{Base-Emitter \(V_{BE}\) and Base-Collector \(V_{BC}\) Voltages}:
        If you rearranged Ebers-Moll model for \(V_{BE}\) or \(V_{BC}\) after
        simplifying for forward active or reverse active mode, you get:
        \begin{itemize}
          \item \textbf{Under Forward Active Mode}:
            \begin{align*}
              V_{BE} = V_T \ln \left( \frac{I_C}{I_{S}} + 1 \right) \\
              V_{BC} \ll 0
            \end{align*}
          \item \textbf{Under Reverse Active Mode}:
            \begin{align*}
              V_{BE} \ll 0 \\
              V_{BC} = V_T \ln \left( \frac{I_C}{I_{S}} + 1 \right)
            \end{align*}
          \end{itemize}
      \item \textbf{Collector-Emitter \(V_{CE}\) Voltage}:
        The collector-emitter voltage is the voltage across the collector and
        emitter terminals. It is given by:
        \begin{align*}
          V_{CE} &= V_{BE} - V_{BC} \\
          V_{CE} &= V_B - V_E - (V_B - V_C) \\
          V_{CE} &= V_B - V_E - V_B + V_C \\
          V_{CE} &= V_C - V_E
        \end{align*}
    \end{itemize}
  \item \textbf{Modes of Operation}:
    A BJT can operate in four different modes depending on \(V_{BE}\) and
    \(V_{CE}\). If \(V_{Fth} \approx 0.7V\), \(V_{Rth} \approx -0.7 V\)
    and \(\beta_R\) is the reverse current gain, then:
    \begin{table}[H]
      \centering
      \begin{tabular}{|c|c|c|c|}\hline
        \textbf{Mode}              & \( V_{BE} \)        & \( V_{BC} \)             & \(I_{C}\) \\ \hline
        Cutoff                     & $V_{BE} < 0.7       $ & $V_{BC} < -0.7 $       & $I_C \approx 0$ \\ \hline
        Forward Active (Amplifier) & $V_{BE} \gtrsim 0.7 $ & $V_{BC} < -0.7 $       & $I_C \approx \beta I_B$ \\ \hline
        Saturation (Switch)        & $V_{BE} \gtrsim 0.7 $ & $V_{BC} \gtrsim -0.7 $ & $I_C \approx I_{C(sat)}$ \\ \hline
        Reverse Active             & $V_{BE} < 0.7       $ & $V_{BC} \gtrsim -0.7 $ & $I_C \approx \beta_R I_B$ \\ \hline
      \end{tabular}
    \end{table}
  \item \textbf{Applications}:
    Various modes can be exploited for different applications. The two
    most common are switching and amplification.
    \begin{itemize}
      \item \textbf{Switching}:
        A BJT can be used as a switch by being driven between \textbf{cutoff}
        (off) and \textbf{saturation} (on). The idea is size the base resistor
        such that there is enough \(I_C\) -- it collapses  \(V_{CE}\) to
        \(V_{CE(sat)}\) (typically \(0.2\text{--}0.3 V\)) through whatever load
        that exists on the collector. You only need to,
        \begin {itemize}
          \item Find the required collector current \(I_C\) to drive the load.
            \[
              I_C = \frac{V_{CC} - V_{CE(sat)}}{R_C}
            \]
          \item Find the required base current \(I_B\) to drive the transistor into
            saturation. \(\beta_F\) is forced current gain to ensure large base current
            and even larger collector current, usually \(\beta_F = 10\).
            \[
              I_B = \frac{I_C}{\beta_{F}}
            \]
          \item Find the base resistor \(R_B\) to limit the base current to \(I_B\).
            \[
              R_B = \frac{V_{IN} - V_{BE}}{I_B}
            \]
        \end{itemize}
      \item \textbf{Amplification}:
        A BJT can be used as an amplifier by being biased in \textbf{forward
        active} mode. The idea is to size the base resistor such that your
        desired \(I_C\) keeps the transistor in forward active mode for the
        entire swing of the input signal. You only need to,
        \begin{itemize}
          \item Find the required \(I_B\) for desired \(I_C\). The forward
            current gain is usually \(\beta = 100\).
            \[
              I_B = \frac{I_C}{\beta}
            \]
          \item Find the base resistor \(R_B\) to limit the base current
            to \(I_B\).
            \[
              R_B = \frac{V_{IN} - V_{BE}}{I_B}
            \]
        \end{itemize}
    \end{itemize}
  \item \textbf{Effects}:
    BJTs are affected by many non-linear effects that are not captured by
    the Ebers-Moll model. Sometimes they cause hard to diagnose problems.
    \begin{itemize}
        \item \textbf{Base Width Modulation}: Increasing \(V_{CE}\) shrinks the base
          width physically. This increases the collector current and is captured
          by the Early voltage (\(V_A\)) and produces the Early Effect.
          \[
            I_C = I_{S} \left( e^{\frac{V_{BE}}{V_T}} - 1 \right) \left( 1 + \frac{V_{CE}}{V_A} \right) \hspace{1cm} \text{[forward active]}
          \]
    \end{itemize}
  \item \textbf{Noise}:
      BJTs are also affected by noise namely, thermal noise, shot noise and flicker
      noise. You can read more about these noises in~\ref{sec:theory_general_noise}.
  \item \textbf{Circuit Analysis}:
    Depending upon Axes of Linearity requirements, You can either use the
    Non-Linear or Linear model to analyse the behaviour of a BJT while
    performing your circuit analysis.
    \begin{itemize}
      \item \textbf{Non-Linear Analysis}:
        Use Ebers-Moll model to analyse the behaviour of a BJT in a circuit.
      \item \textbf{Linear Analysis}:
        Use small-signal model to analyse the behaviour of a BJT in a circuit.
    \end{itemize}
\end{itemize}
